\section{Consideraciones}
\label{sec:consideraciones}

\subsection{Problema y alcance}

El sistema propuesto analiza conjuntos de artículos de Wikipedia en inglés y
francés para obtener información sobre sus fechas, secciones, imágenes y
vocabulario. La entrada proviene de un dataset en formato Parquet disponible
en Kaggle. Cada artículo contiene datos estructurados y un \emph{abstract},
es decir, el texto introductorio que resume su contenido.

El problema consiste en procesar un corpus potencialmente mayor que la memoria
disponible y atender a varios clientes concurrentemente. El diseño distribuye
el trabajo entre computadoras y procesa los datos incrementalmente, preservando
la exactitud de las agregaciones y el aislamiento de cada ejecución.
Agregar recursos debe permitir aumentar tanto la cantidad de clientes
atendidos como el volumen de entrada por cliente.

El cliente dispone de los artículos y los envía incrementalmente al sistema;
la ingesta, la selección de campos y el cálculo de las consultas quedan a cargo
del servicio. Las cinco consultas comparten la entrada de una ejecución.
El cliente también recibe resultados durante el envío y comunica explícitamente
el fin de la entrada. El alcance admite una única ejecución por cliente;
cada ejecución tiene una identidad que permite separar sus datos, estado y
resultados de los de otros clientes.

\subsection{Resultados requeridos}

Para los artículos recibidos en cada ejecución se deben resolver las siguientes
consultas. Sus actores y escenarios se desarrollan en la
sección~\ref{sec:casos-de-uso}.

\begin{description}
    \item[Q1.] URL y fecha de modificación de los artículos en inglés cuya
    fecha de última modificación corresponde a 2020 o a un año anterior.

    \item[Q2.] Nombre, URL del artículo y bytes de su imagen para los artículos
    hasta 2020 inclusive que contienen una sección denominada
    \texttt{biography} o \texttt{biographie}. La búsqueda comprende ambos
    idiomas y compara el nombre completo de la sección en minúsculas.
    El sistema obtiene el contenido de las imágenes referenciadas mediante
    Wikimedia.

    \item[Q3.] Palabras únicas compartidas entre los abstracts de ambos idiomas:
    cada palabra debe aparecer en al menos un artículo en inglés y en al menos
    uno en francés.

    \item[Q4.] Hasta 20 pares de palabra y cantidad, correspondientes a las
    palabras con mayor conteo entre todos los abstracts de la ejecución.

    \item[Q5.] URL y cantidad de palabras del abstract de los artículos desde
    2021 cuyo conteo supera estrictamente el promedio de los abstracts de los
    artículos hasta 2020 inclusive.
\end{description}

Los cortes temporales utilizan la fecha de última modificación
(\texttt{date\_modified}). Para el análisis de vocabulario se normalizan las
palabras a minúsculas y se excluyen tokens numéricos y \emph{stopwords},
palabras triviales de las listas de inglés y francés utilizadas por el sistema.
En Q4 y Q5 cada palabra distinta se cuenta una sola vez por abstract; esta
regla no elimina registros repetidos de la entrada. Por ejemplo, dos apariciones
de una misma palabra válida en un abstract aportan una unidad a Q4 y al conteo
de ese abstract en Q5.

Q1, Q2 y Q3 permiten entregar resultados parciales. Q4 y Q5 requieren haber
procesado toda la entrada del cliente: el top global y el promedio histórico
dependen del conjunto completo. Las respuestas pueden transmitirse en lotes
acotados, acompañados de una finalización explícita de cada consulta, incluso
cuando el resultado sea vacío.

\subsection{Requisitos que condicionan el diseño}

\begin{itemize}
    \item \textbf{Concurrencia y aislamiento.} Cada ejecución conserva entradas,
    estado y resultados aislados al atender varios clientes concurrentemente.
    No se exige paralelismo.

    \item \textbf{Determinismo.} Entradas iguales, bajo el mismo contrato de
    procesamiento, deben producir los mismos resultados aunque cambien el
    orden de llegada o la intercalación entre clientes. En Q4 se ordena por
    cantidad descendente y se desempata por palabra normalizada en orden
    lexicográfico ascendente.

    \item \textbf{Escalabilidad entre computadoras.} El trabajo debe poder
    distribuirse entre más elementos de cómputo. No se exige elasticidad durante
    una ejecución; la propuesta fija las particiones y sus propietarios antes
    de iniciarla, como se explica en la
    sección~\ref{subsec:particionado-aislamiento}.

    \item \textbf{Recursos acotados.} El procesamiento debe limitar memoria,
    lotes y operaciones pendientes. Se permiten archivos temporales locales
    para conservar estado que exceda la RAM y se prohíbe utilizar una base de
    datos.

    \item \textbf{Tecnología y comunicación.} La implementación se realizará
    en Go y utilizará un middleware sobre RabbitMQ que abstraiga la comunicación
    basada en grupos mediante colas y exchanges. La propuesta utiliza TCP entre
    el cliente y Gateway, AMQP para los mensajes internos y HTTPS para obtener
    metadatos y bytes de imágenes de Wikimedia.
\end{itemize}

\subsection{Finalización y fallos}

El fin de entrada del cliente indica que no enviará más artículos. Cada etapa
debe completar su trabajo pendiente antes de propagar la finalización a sus
destinos. Recibir el EOF de un productor no implica haber recibido el de todos
los productores esperados ni haber procesado todos sus datos. Esta distinción
permite que las agregaciones se cierren con resultados completos.

Se requiere un apagado ordenado ante SIGTERM: coordinar el trabajo pendiente,
las confirmaciones y el cierre de archivos y conexiones, o informar una
cancelación explícita si no se puede completar la ejecución. Ante una caída
abrupta, el alcance confirmado es abortar las ejecuciones afectadas. Un nuevo
intento utiliza otra identidad de ejecución y retransmite toda la entrada;
reiniciar un contenedor no recupera por sí solo el trabajo anterior.
La propuesta no ofrece recuperación transparente ni alta disponibilidad.

\subsection{Validación y alcance del documento}

La corrección se verificará mediante un notebook serial de referencia,
adaptado a las mismas listas de stopwords y reglas de conteo del sistema,
incluida la deduplicación por abstract de Q5 y el desempate de Q4.
Sobre un subconjunto fijo del dataset se compararán los resultados esperados
con ejecuciones de un solo cliente y de varios clientes concurrentes.
Variar el orden de entrada permitirá comprobar que el resultado no depende
del orden de procesamiento; las pruebas de Q2 deben utilizar imágenes de
referencia reproducibles para comparar también sus bytes.

Este informe presenta una propuesta de arquitectura mediante vistas $4+1$ y
diagramas de responsabilidades, interacción y despliegue.

